\subsection{An evaluated ordinary-unit obstruction and the corrected three-unit norm coset}

\begin{statement}{E.97}{Proved}{The exact nonzero norm defect of \(\xi c_1\)}
Use the literal roots and fields of E.86:
\[
 \zeta^2+\zeta+1=0,\quad \xi^3=\zeta,\quad
 \beta_1^3=17\zeta,\quad \beta_2^3=\beta_1,\qquad
 L_i=K(\beta_i),\quad B_i=L_i(\xi),\quad K=\mathbb Q(\zeta).
 \tag{E97.1}
\]
Put \(c_1=\beta_1+1+3\zeta\), let \(S\) consist of the primes
above \(3,17\), and write \(E_i=\mathcal O_{L_i}^{\times}\),
\(U_i=\mathcal O_{B_i,S}^{\times}\).
The cyclic generator is always
\(\sigma(\beta_2)=\zeta\beta_2\), fixing \(B_1\).
Its restriction generates \(L_2/L_1\).

In the actual ideal-class coordinates \((a_2,b_2)\) of E86.4
and E86.8, the ordinary-unit norm defect is
\[
 \begin{gathered}
 E_1/N_{L_2/L_1}E_2\xrightarrow{\ \omega_L\ }
 (C_9a_2\oplus C_3b_2)/\langle3a_2\rangle
       \xrightarrow{\sim} C_3^2,\\
 \omega_L(\zeta)=(0,1).
 \end{gathered}
 \tag{E97.2}
\]
The first arrow is an isomorphism. Its literal definition is:
choose \(w\in L_2^\times\) of norm \(u\), solve
\(\operatorname{div}(w)=(\sigma-1)D\), and take
\([D]_3\bmod\langle3a_2\rangle\).
In particular \(\zeta\) is a field norm and has no ordinary-unit
norm preimage.

The corrected field-norm target from E.84 has the exact
S-unit defect
\[
 \begin{gathered}
 \xi c_1\in U_1\cap N_{B_2/B_1}B_2^\times,\qquad
 \omega_B(\xi c_1)=1
      \quad\hbox{in }\operatorname{Cl}_S(B_2)\{3\}=C_3,\\
 \xi c_1\notin N_{B_2/B_1}U_2.
 \end{gathered}
 \tag{E97.3}
\]
Here \(\omega_B\) is exactly the divisor map of E86.14,
and \(1\) is the specified class of the prime
\((73,\xi-55,\beta_2-63)\).
Thus (E97.3) evaluates the remaining coset in E84.14.
For every generator \(\gamma\) of the proved principal ideal
\(Q_1Q_2Q_3\subset\mathcal O_{B_2}\), the ordinary \(B_1\)-unit
\(\xi c_1/N\gamma\) has this same nonzero S-unit norm defect.

All these nonmembership assertions persist on completion
for multiplication by \(3^m\). There is a full calculation on
the resulting one-parameter line:
\[
 \begin{gathered}
 \left\{x\in\widehat U_2:Nx=(\xi c_1)^r\right\}
 =
 \begin{cases}
 \varnothing,&r\notin3\mathbb Z_3,\\
 (\xi c_1)^z\,\ker(N:\widehat U_2\to\widehat U_1),
       &r=3z,
 \end{cases}\\
 \ker N\simeq C_3\oplus\mathbb Z_3^{18}.
 \end{gathered}
 \tag{E97.4}
\]
In the second row \(\xi c_1\) is included from \(B_1\) into
\(B_2\), so its norm is its cube. The kernel in (E97.4)
consists of three-units: its seventeen-valuations are zero.
The abstract direct sum records its compact group type;
the displayed coset uses the actual field norm.

The proof supplies a literal field-norm witness for \(\zeta\),
its complete four-orbit divisor, and its arithmetic Artin
coordinates. It evaluates this finite norm coset without
assigning a value to the eventual exponent \(h\) of E.69
or to the transvection \(W\) of E.91. The comparison to the
positional section coordinates remains the actual norm map
of E.46 and E.88. The manuscript anchor is
LM \(\S51.16\), folios733v--736v, scans172--178.
\end{statement}

\begin{proof}
\emph{The ordinary divisor map and its exact ambiguity.}
E.86 proves
\(\operatorname{Cl}(L_2)\{3\}=C_9a_2\oplus C_3b_2\),
that \(\sigma\) acts trivially on this group, and that
\([\mathfrak q_{3,2}]=3a_2\).
The other ramified prime is \((\beta_2)\) and is principal.
Norm on the ordinary class groups is onto by E.91.
The identity \(jN=1+\sigma+\sigma^2=3\) therefore gives
\[
 j\operatorname{Cl}(L_1)\{3\}=3\operatorname{Cl}(L_2)\{3\}
              =\langle3a_2\rangle .
 \tag{E97.5}
\]
An invariant fractional ideal is generated by extended base
ideals and the ramified primes. To see this on each prime
orbit, its coefficients must all be equal; the orbit sum
is an extended ideal at an unramified prime, while at a
totally ramified prime the fixed upper prime generates
the quotient by the extended ideal. Consequently the
three-primary strongly ambiguous subgroup is exactly
\(\langle3a_2\rangle\).

For completeness, the map in (E97.2) is the inverse of
the exact map in Lemmermeyer, Proposition1, whose full
proof is read here \cite[Proposition1]{E_LemmermeyerAmbiguous}.
Its orientation can be checked directly. A norm-zero
divisor is in the image of \(\sigma-1\), separately on
each finite prime orbit. Thus \(Nw=u\in E_1\) gives
the indicated \(D\). Replacing \(D\) by another solution
adds an invariant ideal. Replacing \(w\) by a second
norm preimage multiplies it by \(\sigma(v)/v\), by
Hilbert90 \cite[II, Proposition1.22 and Corollary1.23]{E_MilneCFT},
and changes \(D\) by a principal ideal, followed by
that invariant ambiguity. Conversely, a zero class modulo
invariant ideals permits \(D\) to be adjusted to a
principal divisor; dividing \(w\) by the corresponding
Hilbert90 element gives an ordinary unit of the same norm.
Every ambiguous ideal class gives a norm of a generator
of its \(\sigma-1\) divisor. This proves both directions
of the isomorphism and its exact kernel.

In this extension every ordinary \(L_1\)-unit is a field
norm. The certified basis in E.64 consists of two free
units and \(\zeta\); their residues at the unique prime
above seventeen are respectively \(\zeta^2,\zeta^2,\zeta\).
All are cubes in \(\mathbb F_{289}^{\times}\).
Thus they are local norms there. At unramified places
ordinary units are local norms, and all infinite places
are complex. The global reciprocity product then kills
the only remaining local norm character, at three.
The cyclic Hasse norm theorem
\cite[VIII, Theorem3.1]{E_MilneCFT} gives the field-norm
assertion. This verifies the domain \(E_1\) in (E97.2).
The quotient has order nine and is killed by three.

\emph{An explicit cyclic-algebra switch.}
Set \(b=\beta_1\), so \(b^6+17b^3+289=0\) and
\(\zeta=b^3/17\). Define \(u_j\in L_1\) by the rows
of this integer matrix:
\[
 \begin{pmatrix}u_0\\u_1\\u_2\end{pmatrix}
 ={1\over11169}
 \begin{pmatrix}
 -717026&-1721029&-474164&29437&-20479&-31853\\
 -2831945&130475&550222&53428&-19498&-6494\\
 -2266763&678062&13447&-275174&-49207&9805
 \end{pmatrix}
 \begin{pmatrix}1\\b\\b^2\\b^3\\b^4\\b^5\end{pmatrix}.
 \tag{E97.6}
\]
Direct reduction modulo the displayed degree-six
polynomial proves the exact identity
\[
 u_0^3+\zeta u_1^3+\zeta^2u_2^3
       -3\zeta u_0u_1u_2=b.
 \tag{E97.7}
\]
Thus \(u_0+u_1\xi+u_2\xi^2\) is a literal norm
preimage of \(b\) in \(B_1/L_1\). The exploratory
computation obtained it from the already certified
\(B_1\) S-units; (E97.6)--(E97.7) verify it independently
of that discovery.

The switch to the required extension is given by actual
matrices over \(L_1\):
\[
 A=\begin{pmatrix}0&0&\zeta\\1&0&0\\0&1&0\end{pmatrix},
 \qquad
 B=(u_0I+u_1A+u_2A^2)
       \operatorname{diag}(1,\zeta,\zeta^2).
 \quad
 A^3=\zeta I,\ B^3=bI,\ BA=\zeta AB.
 \tag{E97.8}
\]
In \(L_2\) take the unique vector \(v\) with
\((B-\beta_2I)v=0\) and \(v_3=1\).
The certificate checks that this eigenspace is
one-dimensional and that this normalization is defined.
Both \(Av\) and \(\sigma(v)\) lie in the
\(\zeta\beta_2\)-eigenline, hence
\[
       Av=w\,\sigma(v),\qquad
       w={ (Av)_1\over \sigma(v)_1}\in L_2^\times .
 \tag{E97.9}
\]
This is a literal definition of the retained witness.
Applying \(A\) three times and using that it commutes
with coefficient-wise \(\sigma\) gives
\(\zeta v=w\,\sigma(w)\sigma^2(w)v\).
Consequently \(N_{L_2/L_1}w=\zeta\).
The program checks this both in the relative field
\(L_1[T]/(T^3-b)\) and in the absolute field
\(\mathbb Q[T]/(T^{18}+17T^9+289)\).

\emph{The full divisor and its actual class.}
For each row of the following table let
\(P=(p,\beta_2-r)\), and order its three primes
as \(P,\sigma P,\sigma^2P\).
The vector \(v_P\) gives their valuations of \(w\);
\(d_P\) gives their coefficients in \(D\).
The last column is the arithmetic Artin value of
the contracted prime
\((p,\beta_1-r^3)\) in the lower genus coordinates
\(\Gamma_L=C_9\oplus C_3^2\) of E86.8:
\[
 \begin{gathered}
 \begin{array}{c|r|r}
 i&p_i&r_i\\ \hline
 1&7&3\\
 2&13&4\\
 3&988728187&449890749\\
 4&212464545716738508121&96597925880218313619
 \end{array}\\[3pt]
 \begin{array}{c|c|c|c}
 i&v_{P_i}&d_{P_i}&\operatorname{Art}(P_i\cap L_1)\\ \hline
 1&(-1,0,1)&(0,0,-1)&(2,1,0)\\
 2&(1,0,-1)&(0,0,1)&(4,2,0)\\
 3&(-1,1,0)&(0,-1,-1)&(1,0,2)\\
 4&(-1,1,0)&(0,-1,-1)&(0,0,0)
 \end{array}
 \end{gathered}
 \tag{E97.10}
\]
All four rational integers \(p\) are proved prime in the
certificate. The prime ideals are kernels of the actual
residue maps on the maximal orders. Their residue degrees
are one, their relative orbits have length three, and
these rows are the complete factorization of \((w)\).
For each row the coefficient identity is
\[
   (d_2-d_0,d_0-d_1,d_1-d_2)=v_P .
 \tag{E97.11}
\]
Multiplication of the full HNF ideals independently
checks \(\sigma(D)/D=(w)\).

The genus field is abelian over \(L_1\); therefore the
three conjugate primes in a row have the same Artin
value in its subgroup over \(L_2\). Their relative
residue degrees are one, so this is exactly the
Artin value of the contracted lower prime. Summing
the last column with coefficients
\(\sum_i d_{P,i}=-1,1,-2,-2\) gives
\[
 \operatorname{Art}(D)=(0,1,2)\in\Gamma_L,\qquad
 (x,y,z)\longmapsto(x-3y-3z,y)
       \quad\Longrightarrow\quad [D]_3=(0,1).
 \tag{E97.12}
\]
This uses the actual E86.8 class map. Its integral
kernel coordinates, including the redundant order-one
coordinate, are checked by the inverse matrices
\[
 C=\begin{pmatrix}1&-3&-3\\0&1&0\\1/3&-2/3&-2/3\end{pmatrix},
 \qquad
 C^{-1}=\begin{pmatrix}-2&0&9\\0&1&0\\-1&-1&3\end{pmatrix}.
 \tag{E97.13}
\]
The source of \(C\) is \(\ker(2(x+y+z)\bmod3)\),
so its last coordinate is integral. Its target orders
are \((9,3,1)\). Quotienting by the ramified class
\((3,0,0)\) gives relation lattice
\(\operatorname{diag}(3,3,1)\); the class of \(D\)
has lattice coordinate \((0,1/3,0)\).
This proves the nonzero value in (E97.2).

\emph{The horizontal transfer of the obstruction.}
On \(B_2\) let
\[
 h(\xi)=\zeta\xi,\qquad h(\beta_2)=\beta_2 .
 \tag{E97.14}
\]
It generates \(B_2/L_2\), and commutes with \(\sigma\).
Thus the actual norms form the commutative square
\[
 \begin{array}{ccc}
 B_2^\times&\xrightarrow{\,N_\sigma\,}&B_1^\times\\
 \big\downarrow N_h&&\big\downarrow N_h\\
 L_2^\times&\xrightarrow{\,N_\sigma\,}&L_1^\times .
 \end{array}
 \tag{E97.15}
\]
E.84 proves the field-norm membership of \(\xi c_1\).
Choose a field preimage \(x\). The unique vertical
prime extensions over each of the three \(Q_i\)
and three \(T_i\) have residue degree one. Since
\((c_1)=\mathfrak q_{3,1}\) in \(L_1\), the six
valuations of \(x\) are \((1,1,1,0,0,0)\).
This statement concerns every field preimage, even
when it has additional support outside S.
The horizontal extension is S-split. Consequently
\[
 y={N_hx\over c_1}\in L_2^\times,\qquad
 v_{\mathfrak q_{3,2}}(y)=v_{\mathfrak q_{17,2}}(y)=0,
 \qquad
 N_\sigma y={N_h(\xi c_1)\over c_1^3}
             ={\zeta c_1^3\over c_1^3}=\zeta .
 \tag{E97.16}
\]
An S-unit \(x\) would make \(y\) an ordinary unit,
contrary to (E97.2). This already proves
the nonmembership in (E97.3).

The same square fixes its exact sign and S-class
coordinate. In the free group of divisors outside S,
write \(\operatorname{div}_S(x)=(\sigma-1)D_B\).
Then
\[
 \operatorname{div}(y)=(\sigma-1)N_hD_B,\qquad
 N_{B_2/L_2}:\operatorname{Cl}_S(B_2)\{3\}\longrightarrow
       \operatorname{Cl}_S(L_2)\{3\},\quad
 z\longmapsto(0,z).
 \tag{E97.17}
\]
The first equality is an equality of full divisors,
because \(y\) has zero S-valuations. The second is
the literal E86.5 map in the coherent prime bases.
The calculation (E97.12) says that its image on
\([D_B]_3=\omega_B(\xi c_1)\) is \((0,1)\).
Thus \(\omega_B(\xi c_1)=1\), as claimed.

E.86 proves that \(Q_1Q_2Q_3\) is principal. A generator
\(\gamma\) has the same six valuations as \(x\).
Accordingly \(\xi c_1/N_\sigma\gamma\) is an ordinary
\(B_1\)-unit. Multiplying by this actual S-unit norm
does not change the class in E97.3, proving the
assertion for every such generator.

\emph{Completion and all powers in this line.}
The unit and S-unit groups here are finitely generated.
Every norm image under consideration contains the cubes
of its target group, through inclusion of a base element.
Their finite cokernels are therefore killed by three;
completion for \(3^m\)-multiplication preserves those
cokernels and the displayed obstruction homomorphisms.
Their completed norm images are also closed, being
images of compact finitely generated \(\mathbb Z_3\)-modules.
Thus the value \(1\in C_3\) in E97.3 is still nonzero
after completion. Its restriction to the powers of
\(\xi c_1\) is \(r\mapsto r\bmod3\), whose kernel is
exactly \(3\mathbb Z_3\).

For \(r=3z\), inclusion of the element
\((\xi c_1)^z\) gives the stated preimage; multiplication
by the actual norm kernel gives every other preimage.
The norm on the six S-valuations has coefficient one
at each prime. Hence a norm-one S-unit has zero
seventeen-valuations. The kernel has free rank
\((27-1+6)-(9-1+6)=18\), and its torsion is exactly
the kernel of cubing on \(\mu_9\), namely \(\mu_3\).
This also follows from the actual kernel calculation
in E.86 and proves (E97.4).

The portable certificate rebuilds only the degree-six
lower genus computation and then verifies (E97.6)--(E97.17).
It retains every field element, eigenvector, ideal HNF,
prime orbit, Artin coordinate, and inverse class map.
The two earlier degree-eighteen BNF attempts remain
separately recorded as running attempts; no upper BNF
or unproved unit basis is used in this proof.
\end{proof}
